\nonstopmode{}
\documentclass[a4paper]{book}
\usepackage[times,inconsolata,hyper]{Rd}
\usepackage{makeidx}
\makeatletter\@ifl@t@r\fmtversion{2018/04/01}{}{\usepackage[utf8]{inputenc}}\makeatother
% \usepackage{graphicx} % @USE GRAPHICX@
\makeindex{}
\begin{document}
\chapter*{}
\begin{center}
{\textbf{\huge Package `BayesianLasso'}}
\par\bigskip{\large \today}
\end{center}
\ifthenelse{\boolean{Rd@use@hyper}}{\hypersetup{pdftitle = {BayesianLasso: Bayesian Lasso Regression and Tools for the Lasso Distribution}}}{}
\ifthenelse{\boolean{Rd@use@hyper}}{\hypersetup{pdfauthor = {John Ormerod; Mohammad Javad Davoudabadi; Garth Tarr; Samuel Mueller; Jonathon Tidswell}}}{}
\begin{description}
\raggedright{}
\item[Type]\AsIs{Package}
\item[Title]\AsIs{Bayesian Lasso Regression and Tools for the Lasso Distribution}
\item[Version]\AsIs{0.3.7}
\item[Date]\AsIs{2026-03-13}
\item[Maintainer]\AsIs{Mohammad Javad Davoudabadi }\email{mohammad.davoudabadi@qut.edu.au}\AsIs{}
\item[Description]\AsIs{Implements Bayesian Lasso regression using efficient Gibbs sampling
algorithms, including modified versions of the Hans and Park–Casella (PC) samplers.
Includes functions for working with the Lasso distribution, such as its density,
cumulative distribution, quantile, and random generation functions, along with moment
calculations. Also includes a function to compute the Mills ratio. Designed for sparse
linear models and suitable for high-dimensional regression problems.}
\item[License]\AsIs{GPL-3}
\item[Imports]\AsIs{Rcpp (>= 1.0.12)}
\item[LinkingTo]\AsIs{Rcpp, RcppNumerical, RcppArmadillo, RcppEigen}
\item[RoxygenNote]\AsIs{7.3.2}
\item[SystemRequirements]\AsIs{C++11}
\item[Encoding]\AsIs{UTF-8}
\item[URL]\AsIs{}\url{https://garthtarr.github.io/BayesianLasso/}\AsIs{,
}\url{https://github.com/garthtarr/BayesianLasso}\AsIs{}
\item[VignetteBuilder]\AsIs{knitr}
\item[Suggests]\AsIs{knitr, rmarkdown, monomvn, posterior, rstan, bayesreg, lars,
Ecdat, testthat (>= 3.0.0), MASS}
\item[Config/testthat/edition]\AsIs{3}
\item[BugReports]\AsIs{}\url{https://github.com/garthtarr/BayesianLasso/issues}\AsIs{}
\item[NeedsCompilation]\AsIs{yes}
\item[Author]\AsIs{John Ormerod [aut, cph] (<}\url{https://orcid.org/0000-0002-4650-7507}\AsIs{>),
Mohammad Javad Davoudabadi [aut, cre, cph]
  (<}\url{https://orcid.org/0000-0001-7312-1530}\AsIs{>),
Garth Tarr [aut, cph] (<}\url{https://orcid.org/0000-0002-6605-7478}\AsIs{>),
Samuel Mueller [aut, cph] (<}\url{https://orcid.org/0000-0002-3087-8127}\AsIs{>),
Jonathon Tidswell [aut, cph]}
\item[Archs]\AsIs{x64}
\end{description}
\Rdcontents{Contents}
\HeaderA{LassoDistribution}{The Lasso Distribution}{LassoDistribution}
\aliasA{dlasso}{LassoDistribution}{dlasso}
\aliasA{elasso}{LassoDistribution}{elasso}
\aliasA{MillsRatio}{LassoDistribution}{MillsRatio}
\aliasA{mlasso}{LassoDistribution}{mlasso}
\aliasA{Modified\_Hans\_Gibbs}{LassoDistribution}{Modified.Rul.Hans.Rul.Gibbs}
\aliasA{Modified\_PC\_Gibbs}{LassoDistribution}{Modified.Rul.PC.Rul.Gibbs}
\aliasA{plasso}{LassoDistribution}{plasso}
\aliasA{qlasso}{LassoDistribution}{qlasso}
\aliasA{rlasso}{LassoDistribution}{rlasso}
\aliasA{vlasso}{LassoDistribution}{vlasso}
\aliasA{zlasso}{LassoDistribution}{zlasso}
%
\begin{Description}
Provides functions related to the Lasso distribution, including the normalizing constant,
probability density function, cumulative distribution function, quantile function, and
random number generation for given parameters \code{a}, \code{b}, and \code{c}.
Additional utilities include the Mills ratio, expected value, and variance of the distribution.
The package also implements modified versions of the Hans and Park–Casella Gibbs sampling algorithms
for Bayesian Lasso regression.
\end{Description}
%
\begin{Usage}
\begin{verbatim}
zlasso(a, b, c, logarithm)
dlasso(x, a, b, c, logarithm)
plasso(q, a, b, c)
qlasso(p, a, b, c)
rlasso(n, a, b, c)
elasso(a, b, c)
vlasso(a, b, c)
mlasso(a, b, c)
MillsRatio(d)
Modified_Hans_Gibbs(X, y, beta_init, a1, b1, u1, v1,
              nsamples, lambda_init, sigma2_init, thin, verbose,
              tune_lambda2, rao_blackwellization)
Modified_PC_Gibbs(X, y, a1, b1, u1, v1, 
              nsamples, lambda_init, sigma2_init, thin, verbose)
\end{verbatim}
\end{Usage}
%
\begin{Arguments}
\begin{ldescription}
\item[\code{x}, \code{q}] Vector of quantiles (vectorized).

\item[\code{p}] Vector of probabilities.

\item[\code{a}] Vector of precision parameter which must be non-negative.

\item[\code{b}] Vector of off set parameter.

\item[\code{c}] Vector of tuning parameter which must be non-negative values.

\item[\code{n}] Number of observations.

\item[\code{logarithm}] Logical. If \code{TRUE}, probabilities are returned on the log scale.

\item[\code{d}] A scalar numeric value. Represents the point at which the Mills ratio is evaluated.

\item[\code{X}] Design matrix (numeric matrix).

\item[\code{y}] Response vector (numeric vector).

\item[\code{a1}] Shape parameter of the prior on \eqn{\lambda^2}{}.

\item[\code{b1}] Rate parameter of the prior on \eqn{\lambda^2}{}.

\item[\code{u1}] Shape parameter of the prior on \eqn{\sigma^2}{}.

\item[\code{v1}] Rate parameter of the prior on \eqn{\sigma^2}{}.

\item[\code{nsamples}] Number of Gibbs samples to draw.

\item[\code{beta\_init}] Initial value for the model parameter \eqn{\beta}{}.

\item[\code{lambda\_init}] Initial value for the shrinkage parameter \eqn{\lambda^2}{}.

\item[\code{sigma2\_init}] Initial value for the error variance \eqn{\sigma^2}{}.

\item[\code{thin}] Thinning interval for the MCMC chain. Only every `thin`-th draw is stored. Default is 1 (no thinning).

\item[\code{verbose}] Integer. If greater than 0, progress is printed every \code{verbose} iterations during sampling. Set to 0 to suppress output.

\item[\code{tune\_lambda2}] Logical; if TRUE (default), the tuning parameter
\eqn{\lambda^2}{} is estimated during sampling.

\item[\code{rao\_blackwellization}] Logical; if TRUE, Rao–Blackwellization
is applied to improve posterior estimation. Default is FALSE.
\end{ldescription}
\end{Arguments}
%
\begin{Details}
If \eqn{X \sim \text{Lasso}(a, b, c)}{} then its density function is:
\deqn{
p(x;a,b,c) = Z^{-1} \exp\left(-\frac{1}{2} a x^2 + bx - c|x| \right)
}{}
where \eqn{x \in \mathbb{R}}{}, \eqn{a > 0}{}, \eqn{b \in \mathbb{R}}{}, \eqn{c > 0}{}, and \eqn{Z}{} is the normalizing constant.

More details are included for the CDF, quantile function, and normalizing constant in the original documentation.
\end{Details}
%
\begin{Value}
\begin{itemize}


\item{} \code{zlasso}, \code{dlasso}, \code{plasso}, \code{qlasso}, \code{rlasso}, \code{elasso}, \code{vlasso}, \code{mlasso}, \code{MillsRatio}: 
return the corresponding scalar or vector values related to the Lasso distribution and a numeric value representing the Mills ratio.
\item{} \code{Modified\_Hans\_Gibbs}: returns a list containing:
\begin{description}

\item[\code{mBeta}] Matrix of MCMC samples for the regression coefficients \eqn{\beta}{}, with \code{nsamples} rows and \code{p} columns.
\item[\code{vsigma2}] Vector of MCMC samples for the error variance \eqn{\sigma^2}{}.
\item[\code{vlambda2}] Vector of MCMC samples for the shrinkage parameter \eqn{\lambda^2}{}.
\item[\code{mA}] Matrix of sampled values for parameter \eqn{a_j}{} of the Lasso distribution for each \eqn{\beta_j}{}.
\item[\code{mB}] Matrix of sampled values for parameter \eqn{b_j}{} of the Lasso distribution for each \eqn{\beta_j}{}.
\item[\code{mC}] Matrix of sampled values for parameter \eqn{c_j}{} of the Lasso distribution for each \eqn{\beta_j}{}.

\end{description}

\item{} \code{Modified\_PC\_Gibbs}: returns a list containing:
\begin{description}

\item[\code{mBeta}] Matrix of MCMC samples for the regression coefficients \eqn{\beta}{}.
\item[\code{vsigma2}] Vector of MCMC samples for the error variance \eqn{\sigma^2}{}.
\item[\code{vlambda2}] Vector of MCMC samples for the shrinkage parameter \eqn{\lambda^2}{}.
\item[\code{mM}] Matrix of estimated means of the full conditional distributions of each \eqn{\beta_j}{}.
\item[\code{mV}] Matrix of estimated variances of the full conditional distributions of each \eqn{\beta_j}{}.
\item[\code{va\_til}] Vector of estimated shape parameters for the full conditional inverse-gamma distribution of \eqn{\sigma^2}{}.
\item[\code{vb\_til}] Vector of estimated rate parameters for the full conditional inverse-gamma distribution of \eqn{\sigma^2}{}.
\item[\code{vu\_til}] Vector of estimated shape parameters for the full conditional inverse-gamma distribution of \eqn{\lambda^2}{}.
\item[\code{vv\_til}] Vector of estimated rate parameters for the full conditional inverse-gamma distribution of \eqn{\lambda^2}{}.

\end{description}



\end{itemize}

\end{Value}
%
\begin{SeeAlso}
\code{\LinkA{normalize}{normalize}} for preprocessing input data before applying the samplers.
\end{SeeAlso}
%
\begin{Examples}
\begin{ExampleCode}
a <- 2; b <- 1; c <- 3
x <- seq(-3, 3, length.out = 1000)
plot(x, dlasso(x, a, b, c, logarithm = FALSE), type = 'l')

r <- rlasso(1000, a, b, c)
hist(r, breaks = 50, probability = TRUE, col = "grey", border = "white")
lines(x, dlasso(x, a, b, c, logarithm = FALSE), col = "blue")

plasso(0, a, b, c)
qlasso(0.25, a, b, c)
elasso(a, b, c)
vlasso(a, b, c)
mlasso(a, b, c)
MillsRatio(2)




# The Modified_Hans_Gibbs() function uses the Lasso distribution to draw 
# samples from the full conditional distribution of the regression coefficients.

y <- 1:20
X <- matrix(c(1:20,12:31,7:26),20,3,byrow = TRUE)

a1 <- b1 <- u1 <- v1 <- 0.01
sigma2_init <- 1
lambda_init <- 0.1
beta_init <- rep(1, ncol(X))
nsamples <- 1000
verbose <- 100
tune_lambda2 <- TRUE
rao_blackwellization <- FALSE

Output_Hans <- Modified_Hans_Gibbs(
                X, y, beta_init, a1, b1, u1, v1,
                nsamples, lambda_init, sigma2_init, 
                verbose, tune_lambda2, rao_blackwellization
)

colMeans(Output_Hans$mBeta)
mean(Output_Hans$vlambda2)


Output_PC <- Modified_PC_Gibbs(
               X, y, a1, b1, u1, v1, 
               nsamples, lambda_init, sigma2_init, verbose)

colMeans(Output_PC$mBeta)
mean(Output_PC$vlambda2)

\end{ExampleCode}
\end{Examples}
\HeaderA{normalize}{Normalize Response and Covariates}{normalize}
%
\begin{Description}
This function centers and (optionally) scales the response vector and each column
of the design matrix using the population variance. It is used to prepare data
for Bayesian Lasso regression.
\end{Description}
%
\begin{Usage}
\begin{verbatim}
normalize(y, X, scale = TRUE)
\end{verbatim}
\end{Usage}
%
\begin{Arguments}
\begin{ldescription}
\item[\code{y}] A numeric response vector.

\item[\code{X}] A numeric matrix or data frame of covariates (design matrix).

\item[\code{scale}] Logical; if \code{TRUE}, variables are scaled to have unit population variance (default is \code{TRUE}).
\end{ldescription}
\end{Arguments}
%
\begin{Value}
A list with the following elements:
\begin{itemize}

\item{} \code{vy}: Normalized response vector.
\item{} \code{mX}: Normalized design matrix.
\item{} \code{mu.y}: Mean of the response vector.
\item{} \code{sigma2.y}: Population variance of the response vector.
\item{} \code{mu.x}: Vector of column means of \code{X}.
\item{} \code{sigma2.x}: Vector of population variances for columns of \code{X}.

\end{itemize}

\end{Value}
%
\begin{Examples}
\begin{ExampleCode}
set.seed(1)
X <- matrix(rnorm(100 * 10), 100, 10)
beta <- c(2, -3, rep(0, 8))
y <- as.vector(X %*% beta + rnorm(100))
norm_result <- normalize(y, X)

\end{ExampleCode}
\end{Examples}
\printindex{}
\end{document}
